\documentclass[11pt]{article}
\usepackage[a4paper,margin=2.5cm]{geometry}
\usepackage{amsmath,amssymb,amsthm}
\usepackage{bm}
\usepackage{graphicx}
\usepackage{float}
\graphicspath{{"../Part 1 code/outputs/"}{"../Part 2 code/outputs/"}{"../Part 3 code/outputs/"}{"../Part 4 code/outputs/"}{"../Part 5 code/outputs/"}{"../Part 6 code/outputs/"}}

\newcommand{\E}{\mathbb{E}}
\newcommand{\Var}{\operatorname{Var}}
\newcommand{\N}{\mathcal{N}}
\newcommand{\dd}{\mathrm{d}}

\title{Assignment 1}
\author{Mohd Zeeshan \\ Roll no.\ 2025901020}
\date{}

\begin{document}
\maketitle

\section*{Part 1: Simulating the Forward Diffusion}

\subsection*{Derivations}

Throughout, $X_t$ solves the Ornstein--Uhlenbeck SDE
\begin{equation}\label{eq:ou}
\dd X_t = -\theta X_t\,\dd t + \sqrt{2\theta}\,\dd W_t, \qquad \theta>0,
\end{equation}
with initial density
$p_0(x) = \tfrac12\varphi(x\mid -3,0.5^2) + \tfrac12\varphi(x\mid 3,0.5^2)$.
We use the It\^o rules $(\dd t)^2 = 0$, $\dd t\,\dd W_t = 0$, $(\dd W_t)^2 = \dd t$,
and $\E[\dd W_t]=0$ (since $\dd W_t\sim\N(0,\dd t)$ and is independent of $X_t$).

\subsubsection*{1(a) Stationary distribution}

\paragraph{Mean.} Taking expectations of \eqref{eq:ou},
\[
\dd\,\E[X_t] = -\theta\,\E[X_t]\,\dd t + \sqrt{2\theta}\,\E[\dd W_t]
\quad\Longrightarrow\quad
\frac{\dd m}{\dd t} = -\theta\, m(t).
\]
Separating variables, $\int \frac{\dd m}{m} = -\theta\int \dd t$, so
$\ln|m| = -\theta t + c$, i.e.\ $m(t) = C e^{-\theta t}$. Setting $t=0$ gives $C=m(0)$, hence
\[
m(t) = m(0)\,e^{-\theta t} \xrightarrow[t\to\infty]{} 0 .
\]

\paragraph{Second moment.} Apply It\^o's lemma,
$\dd f = \partial_t f\,\dd t + \partial_x f\,\dd X_t + \tfrac12\partial_{xx} f\,(\dd X_t)^2$,
to $f(x)=x^2$, for which $\partial_t f = 0$, $\partial_x f = 2x$, $\partial_{xx}f = 2$:
\begin{equation}\label{eq:ito-sq}
\dd(X_t^2) = 2X_t\,\dd X_t + (\dd X_t)^2 .
\end{equation}
Squaring \eqref{eq:ou},
\[
(\dd X_t)^2 = \theta^2 X_t^2 (\dd t)^2 + 2\theta(\dd W_t)^2 - 2\theta\sqrt{2\theta}\,X_t\,\dd t\,\dd W_t = 2\theta\,\dd t .
\]
Substituting this and \eqref{eq:ou} into \eqref{eq:ito-sq},
\[
\dd(X_t^2) = 2X_t\bigl(-\theta X_t\,\dd t + \sqrt{2\theta}\,\dd W_t\bigr) + 2\theta\,\dd t
= \bigl(-2\theta X_t^2 + 2\theta\bigr)\dd t + 2\sqrt{2\theta}\,X_t\,\dd W_t .
\]
Taking expectations (the It\^o integral term has zero mean),
\begin{equation}\label{eq:second}
\frac{\dd\,\E[X_t^2]}{\dd t} = -2\theta\,\E[X_t^2] + 2\theta .
\end{equation}

\paragraph{Variance.} By definition $\Sigma(t) = \E[X_t^2] - m(t)^2$, so
\[
\frac{\dd\Sigma}{\dd t} = \frac{\dd\,\E[X_t^2]}{\dd t} - 2m(t)\frac{\dd m}{\dd t}.
\]
Using \eqref{eq:second} and $\dd m/\dd t = -\theta m$,
\[
\frac{\dd\Sigma}{\dd t} = -2\theta\,\E[X_t^2] + 2\theta + 2\theta\, m(t)^2
= -2\theta\bigl(\E[X_t^2]-m(t)^2\bigr) + 2\theta
= -2\theta\,\Sigma(t) + 2\theta .
\]
Separating variables, $\int_0^t \frac{\dd\Sigma}{\Sigma-1} = -2\theta\int_0^t \dd s$, so
$\ln\left|\frac{\Sigma(t)-1}{\Sigma(0)-1}\right| = -2\theta t$ and
\[
\Sigma(t) = 1 + \bigl(\Sigma(0)-1\bigr)e^{-2\theta t} \xrightarrow[t\to\infty]{} 1 .
\]

\paragraph{Stationary law.} The limiting mean and variance are $0$ and $1$. To identify the
whole distribution (not only its first two moments) we use the solution derived in 1(b):
conditionally on $X_0$, $X_t\sim\N(e^{-\theta t}X_0,\,1-e^{-2\theta t})$, whose characteristic function is
$\E[e^{iuX_t}\mid X_0] = \exp\!\bigl(iue^{-\theta t}X_0 - \tfrac12u^2(1-e^{-2\theta t})\bigr)$.
As $t\to\infty$ this tends to $e^{-u^2/2}$ for every $X_0$, and by dominated convergence
$\E[e^{iuX_t}]\to e^{-u^2/2}$ for any initial law. Hence
\[
\boxed{\pi = \N(0,1).}
\]
Equivalently, $\N(0,1)$ is stationary: if $X_0\sim\N(0,1)$ then $X_t$ is Gaussian with
$m(t)=0$ and $\Sigma(t)=1$ for all $t$.

\subsubsection*{1(b) Solution by integrating factor}

Write \eqref{eq:ou} as $\dd X_t + \theta X_t\,\dd t = \sqrt{2\theta}\,\dd W_t$, which has the linear form
$\dot x + P(t)x = Q(t)$ with $P(t)=\theta$. The integrating factor is
$I_t = \exp\!\left(\int\theta\,\dd t\right) = e^{\theta t}$. Define
$Y_t = f(t,X_t) = e^{\theta t}X_t$. The partial derivatives are
\[
\partial_t f = \theta e^{\theta t}x,\qquad
\partial_x f = e^{\theta t},\qquad
\partial_{xx} f = 0,\qquad
\partial_{tx} f = \theta e^{\theta t},\qquad
\partial_{tt} f = \theta^2 e^{\theta t}x .
\]
The multivariate It\^o expansion is
\[
\dd Y_t = \partial_t f\,\dd t + \partial_x f\,\dd X_t + \tfrac12\partial_{xx}f\,(\dd X_t)^2
+ \partial_{tx}f\,\dd t\,\dd X_t + \tfrac12\partial_{tt}f\,(\dd t)^2 .
\]
Here $(\dd t)^2=0$, $\dd t\,\dd X_t = -\theta X_t(\dd t)^2 + \sqrt{2\theta}\,\dd t\,\dd W_t = 0$, and the
$(\dd X_t)^2 = 2\theta\,\dd t$ term is multiplied by $\partial_{xx}f=0$. Therefore
\[
\dd(e^{\theta t}X_t) = \theta e^{\theta t}X_t\,\dd t + e^{\theta t}\bigl(-\theta X_t\,\dd t + \sqrt{2\theta}\,\dd W_t\bigr)
= \sqrt{2\theta}\,e^{\theta t}\,\dd W_t .
\]
Integrating from $0$ to $t$,
\[
e^{\theta t}X_t - X_0 = \sqrt{2\theta}\int_0^t e^{\theta s}\,\dd W_s
\quad\Longrightarrow\quad
X_t = e^{-\theta t}X_0 + \sqrt{2\theta}\int_0^t e^{-\theta(t-s)}\,\dd W_s .
\]

\paragraph{Gaussianity.} The stochastic integral has a deterministic integrand, so it is a limit of
sums $\sum_k g(s_k)(W_{s_{k+1}}-W_{s_k})$ of independent Gaussians and is therefore Gaussian. It is
independent of $X_0$, so $X_t\mid X_0$ is Gaussian and is fixed by its mean and variance.

\paragraph{Mean.} Since $\E\bigl[\int_0^t g(s)\,\dd W_s\bigr]=0$,
\[
\E[X_t\mid X_0] = e^{-\theta t}X_0 + \sqrt{2\theta}\int_0^t e^{-\theta(t-s)}\,\E[\dd W_s] = e^{-\theta t}X_0 .
\]

\paragraph{Variance.} $X_t - \E[X_t\mid X_0] = \sqrt{2\theta}\int_0^t e^{-\theta(t-s)}\,\dd W_s$. By It\^o's isometry,
$\E\bigl[(\int_0^t g(s)\,\dd W_s)^2\bigr] = \int_0^t g(s)^2\,\dd s$, so
\[
\Var(X_t\mid X_0) = 2\theta\int_0^t e^{-2\theta(t-s)}\,\dd s
= 2\theta\, e^{-2\theta t}\left[\frac{e^{2\theta s}}{2\theta}\right]_0^t
= e^{-2\theta t}\bigl(e^{2\theta t}-1\bigr) = 1 - e^{-2\theta t}.
\]
Hence
\[
\boxed{X_t\mid X_0 \sim \N\!\left(e^{-\theta t}X_0,\; 1-e^{-2\theta t}\right).}
\]

\subsubsection*{1(c) Law of $X_t$}

From 1(b) we may write $X_t = e^{-\theta t}X_0 + \xi_t$ with $\xi_t\sim\N(0,1-e^{-2\theta t})$ independent of $X_0$.
Introduce a latent label $Z\in\{1,2\}$ with $P(Z=1)=P(Z=2)=\tfrac12$ such that
\[
X_0\mid Z=1 \sim \N(-3,0.5^2),\qquad X_0\mid Z=2 \sim \N(3,0.5^2),
\]
which reproduces $p_0$. By the law of total probability,
$p_t(x) = \tfrac12\,p(x\mid Z=1) + \tfrac12\,p(x\mid Z=2)$.
Conditionally on $Z$, $X_t$ is a sum of two independent Gaussians, hence Gaussian, with
\[
\E[X_t\mid Z=1] = e^{-\theta t}\E[X_0\mid Z=1] + \E[\xi_t] = -3e^{-\theta t},\qquad
\E[X_t\mid Z=2] = 3e^{-\theta t},
\]
and, using $\Var(cY)=c^2\Var(Y)$ and independence,
\[
\Var(X_t\mid Z=k) = e^{-2\theta t}\Var(X_0\mid Z=k) + \Var(\xi_t)
= 0.25\,e^{-2\theta t} + 1 - e^{-2\theta t} = 1 - 0.75\,e^{-2\theta t} =: v_t ,
\]
for $k=1,2$. Therefore
\[
\boxed{p_t(x) = \tfrac12\varphi\!\left(x\mid -3e^{-\theta t}, v_t\right) + \tfrac12\varphi\!\left(x\mid 3e^{-\theta t}, v_t\right),
\qquad v_t = 1-0.75\,e^{-2\theta t}.}
\]

\paragraph{Checks.} At $t=0$: $v_0 = 1-0.75 = 0.25 = 0.5^2$, and the means are $\mp3$, recovering $p_0$.
As $t\to\infty$: $v_t = 1-0.75\lim_{t\to\infty}e^{-2\theta t} = 1$ and both means tend to $0$, so
$p_t\to\varphi(x\mid0,1)=\pi$, consistent with 1(a).

\subsubsection*{1(d) Dependence on $\theta t$ only}

Set $s=\theta t$. Then
\[
p_t(x) = \tfrac12\varphi\!\left(x\mid -3e^{-s},\,1-0.75e^{-2s}\right) + \tfrac12\varphi\!\left(x\mid 3e^{-s},\,1-0.75e^{-2s}\right).
\]
Neither $\theta$ nor $t$ appears anywhere except through $s$, so the law of $X_t$ depends on $(\theta,t)$
only through the product $\theta t$.

This also follows directly from the SDE. Let $\tilde X_s = X_{s/\theta}$. Since
$\tilde W_s = \sqrt{\theta}\,W_{s/\theta}$ is a standard Brownian motion,
\[
\dd\tilde X_s = -\theta\tilde X_s\,\frac{\dd s}{\theta} + \sqrt{2\theta}\,\frac{\dd\tilde W_s}{\sqrt\theta}
= -\tilde X_s\,\dd s + \sqrt2\,\dd\tilde W_s ,
\]
which is the $\theta=1$ OU process with the same initial law. Changing $\theta$ is therefore only a
rescaling of time.

\subsection*{Implementation}

Equation~\eqref{eq:ou} is discretised with the Euler--Maruyama scheme
\[
X_{t+\Delta t} = X_t - \theta X_t\,\Delta t + \sqrt{2\theta\Delta t}\,Z_t,\qquad Z_t\overset{\text{iid}}{\sim}\N(0,1),
\]
using the defaults $\theta=1$, $N=2\times10^4$, $\Delta t=10^{-3}$ and a fixed random seed.
The horizon is $T=5.40$, the value obtained in Part 2(d). The code and instructions to reproduce every figure
are in the \texttt{Part 1 code} folder.

\subsubsection*{(a) Sample paths}

Figure~\ref{fig:paths} shows ten trajectories (simulated with a separate seed, so that the figure
shows typical paths rather than a rare large excursion). Paths start near $\pm3$, are pulled towards $0$ on a
time scale of order $1/\theta$, and then fluctuate with unit variance, mixing between the two
former modes.

\begin{figure}[H]
\centering
\includegraphics[width=0.8\linewidth]{a_sample_paths.png}
\caption{Ten Euler--Maruyama sample paths of the OU process, $\theta=1$.}
\label{fig:paths}
\end{figure}

\subsubsection*{(b) Empirical histograms versus the exact density}

Figure~\ref{fig:hist} overlays the histogram of the $N$ simulated particles (100 bins on $[-5,5]$)
on the exact density $p_t$ from 1(c) at five equally spaced times in $[0,T]$. We measure the
discrepancy as $\max_j |\hat p_t(x_j) - p_t(x_j)|$ over bin centres $x_j$.

\begin{figure}[H]
\centering
\includegraphics[width=\linewidth]{b_histograms.png}
\caption{Simulated histograms (shaded) against the exact density $p_t$ (black).}
\label{fig:hist}
\end{figure}

\begin{table}[H]
\centering
\begin{tabular}{cccccc}
\hline
$t$ & 0.00 & 1.35 & 2.70 & 4.05 & 5.40 \\
\hline
$\max_j|\hat p_t - p_t|$ & 0.0244 & 0.0149 & 0.0210 & 0.0248 & 0.0295 \\
\hline
\end{tabular}
\caption{Largest pointwise histogram error at each snapshot time.}
\end{table}

The largest discrepancy observed is $0.030$. This error is about what Monte Carlo noise alone
predicts: with bin width $h=0.1$ and $N=2\times10^4$, the standard error of a bin height near a peak
of height $\approx 0.4$ is $\sqrt{p/(Nh)}\approx 0.014$, and the maximum over 100 bins is a few such
standard errors. The error does not grow with $t$, so the time-discretisation bias is negligible
at $\Delta t = 10^{-3}$.

\subsubsection*{(c) Disappearance of bimodality}

Under the dynamics both mixture components move towards the origin, with means
$\pm3e^{-\theta t}$, while their variance $v_t$ grows from $0.25$ to $1$. The two bumps therefore
approach and widen until they merge into a single peak (Figure~\ref{fig:bimodal}).

For a symmetric mixture $\tfrac12\varphi(x\mid -\mu,v) + \tfrac12\varphi(x\mid \mu,v)$, the origin is always a
critical point, and
\[
p''(0) \propto \frac{\mu^2}{v} - 1 .
\]
Hence the density is bimodal (with a local minimum at $0$) iff $\mu^2 > v$, and unimodal otherwise. With
$\mu=3e^{-\theta t}$ and $v_t = 1-0.75e^{-2\theta t}$, the condition $\mu^2 = v_t$ gives
$9.75\,e^{-2\theta t} = 1$, i.e.
\[
t^* = \frac{\ln 9.75}{2\theta} \approx 1.1386 \quad (\theta=1).
\]

\paragraph{Numerical check.} Counting local maxima of the exact $p_t$ on a grid of 4001 points in
$[-6,6]$, stepping $t$ by $10^{-4}$, gives a transition at $t=1.1387$, matching $t^*$.

\begin{figure}[H]
\centering
\includegraphics[width=0.75\linewidth]{c_bimodality.png}
\caption{Exact $p_t$ at several times, including $t^*$.}
\label{fig:bimodal}
\end{figure}

\subsubsection*{(d) Collapse under rescaled time}

We simulate $\theta\in\{0.5,1,2\}$ and take snapshots at $t = s/\theta$ for $s=\theta t\in\{0, 1.35, 2.70, 4.05, 5.40\}$.
Figure~\ref{fig:collapse} shows that the three histograms coincide at each $s$ and match the
exact density evaluated at $\theta t = s$. All pointwise errors lie between $0.015$ and $0.031$,
the same Monte Carlo noise level as in (b). This confirms derivation 1(d): the law of $X_t$
depends on $\theta$ and $t$ only through $\theta t$.

\begin{figure}[H]
\centering
\includegraphics[width=\linewidth]{d_rescaled_time.png}
\caption{Histograms for $\theta\in\{0.5,1,2\}$ at equal rescaled times $\theta t$, with the exact density.}
\label{fig:collapse}
\end{figure}

\section*{Part 2: Convergence to Equilibrium in Wasserstein Distance}

\subsection*{The one-dimensional $W_2$ estimator}

For probability measures $\mu,\nu$ on $\mathbb R$ with finite second moments,
$W_2^2(\mu,\nu) = \inf_{\gamma\in\Gamma(\mu,\nu)} \E_{(X,Y)\sim\gamma}|X-Y|^2$, where $\Gamma(\mu,\nu)$ is the set of couplings.
Given samples $x_1,\dots,x_n\sim\mu$ and $y_1,\dots,y_n\sim\nu$, we use
\[
\widehat W_2^2 = \frac1n\sum_{i=1}^n \bigl(x_{(i)} - y_{(i)}\bigr)^2,
\]
where $x_{(i)}$ and $y_{(i)}$ are the order statistics.

\paragraph{Why sorting is optimal.} Between two sets of $n$ equally weighted points, an optimal coupling can be taken
to be a one-to-one matching of the $x_i$ with the $y_j$. Suppose two matched pairs cross,
i.e.\ $x_a < x_b$ but $y_a > y_b$. Then
\[
(x_a-y_a)^2 + (x_b-y_b)^2 - \bigl[(x_a-y_b)^2 + (x_b-y_a)^2\bigr] = 2(x_b-x_a)(y_a-y_b) > 0,
\]
so swapping the partners strictly lowers the cost. Hence an optimal matching has no crossings, which means it pairs
the $i$-th smallest $x$ with the $i$-th smallest $y$. So $\widehat W_2$ is exactly the $W_2$ distance between the two
sample sets, and it converges to $W_2(\mu,\nu)$ as $n\to\infty$. Because both samples are random, $\widehat W_2$ does not reach
$0$ even for identical distributions; it levels off at a small noise floor that shrinks as $n$ grows.

\subsection*{2(c) Proof of the contraction bound by synchronous coupling}

Let $\gamma^\star$ be an optimal coupling of $p_0$ and $\pi$ (it exists since both have finite second moments), and draw
$(X_0,Y_0)\sim\gamma^\star$, so that $\E|X_0-Y_0|^2 = W_2^2(p_0,\pi)$. Drive two copies of~\eqref{eq:ou} with the
\emph{same} Brownian motion $W_t$:
\[
\dd X_t = -\theta X_t\,\dd t + \sqrt{2\theta}\,\dd W_t,\qquad
\dd Y_t = -\theta Y_t\,\dd t + \sqrt{2\theta}\,\dd W_t .
\]
Subtracting, the noise cancels exactly and $D_t = X_t - Y_t$ satisfies the deterministic ODE
\[
\dd D_t = -\theta D_t\,\dd t \quad\Longrightarrow\quad D_t = e^{-\theta t} D_0 .
\]
Now $X_t\sim p_t$. Since $\pi$ is stationary (1(a)) and $Y_0\sim\pi$, also $Y_t\sim\pi$ for all $t$. Thus $(X_t,Y_t)$ is a
coupling of $p_t$ and $\pi$, and by the definition of $W_2$ as an infimum over couplings,
\[
W_2^2(p_t,\pi) \le \E|X_t-Y_t|^2 = e^{-2\theta t}\,\E|X_0-Y_0|^2 = e^{-2\theta t}\,W_2^2(p_0,\pi).
\]
Taking square roots gives
\[
\boxed{W_2(p_t,\pi) \le e^{-\theta t}\,W_2(p_0,\pi).}
\]

\subsection*{2(d) Choice of the horizon $T$}

The bound guarantees $W_2(p_T,\pi)\le\varepsilon$ as soon as $e^{-\theta T}W_2(p_0,\pi)\le\varepsilon$, so the smallest such time is
\[
T = \frac1\theta\ln\frac{W_2(p_0,\pi)}{\varepsilon}, \qquad \varepsilon = 10^{-2}.
\]
The sorted-sample estimate with $N = 2\times10^4$ gives $W_2(p_0,\pi) = 2.2106$, so with $\theta = 1$ we get
$T = 5.3984$. We round up and use
\[
\boxed{T = 5.40}
\]
for the rest of the assignment, including the Part 1 figures.

\subsection*{Implementation}

All distances below use the sorted-sample estimator with $N=2\times10^4$ samples on each side. Samples of $p_t$ are the
Euler--Maruyama particles from Part 1 ($\theta=1$, $\Delta t=10^{-3}$). They are compared with a fixed set of fresh
samples from $\pi$ and from $p_0$, drawn independently of the simulation. Distances are computed at 55 equally
spaced times in $[0,T]$.

\subsubsection*{2(a) Distance to equilibrium and to the start}

Figure~\ref{fig:w2} shows $W_2(p_t,\pi)$ falling from $2.21$ at $t=0$ to about $0.02$, while $W_2(p_t,p_0)$ rises from
near $0$ to $2.21$: as the process approaches equilibrium it moves away from where it started. The curves
cross at $t\approx0.5$, where $p_t$ is equally far from both. By $t\approx3$ both
curves are flat, so the process has essentially reached $\pi$, and the distance back to $p_0$ equals $W_2(p_0,\pi)$.

Neither curve reaches exactly $0$. At $t=0$, $W_2(p_t,p_0)$ is $0.18$ rather than $0$, and at late times $W_2(p_t,\pi)$ levels
off at about $0.02$. Both are the noise floor of the estimator: two independent sample sets from the same
distribution are still a positive distance apart. The floor is larger for $p_0$ because few samples fall in the
gap between its two bumps, so the sorted values in that region line up poorly.

\begin{figure}[H]
\centering
\includegraphics[width=0.8\linewidth]{a_w2_curves.png}
\caption{Estimated $W_2(p_t,\pi)$ and $W_2(p_t,p_0)$ over $[0,T]$, $\theta=1$.}
\label{fig:w2}
\end{figure}

\subsubsection*{2(b) Contraction bound}

The estimate of the initial distance is $W_2(p_0,\pi) = 2.2106$. Figure~\ref{fig:contraction} compares the bound
$e^{-\theta t}W_2(p_0,\pi)$ with the estimated $W_2(p_t,\pi)$, and Table~\ref{tab:contraction} lists selected values.

\begin{figure}[H]
\centering
\includegraphics[width=0.8\linewidth]{b_contraction.png}
\caption{Contraction bound (dashed) and estimated $W_2(p_t,\pi)$ (dots), $\theta=1$.}
\label{fig:contraction}
\end{figure}

\begin{table}[H]
\centering
\begin{tabular}{cccc}
\hline
$t$ & estimate & bound & bound / estimate \\
\hline
0.0 & 2.208 & 2.211 & 1.00 \\
0.9 & 0.562 & 0.899 & 1.60 \\
1.8 & 0.105 & 0.365 & 3.49 \\
2.7 & 0.028 & 0.149 & 5.31 \\
3.6 & 0.025 & 0.060 & 2.41 \\
4.5 & 0.020 & 0.025 & 1.20 \\
5.4 & 0.017 & 0.010 & 0.58 \\
\hline
\end{tabular}
\caption{Estimated $W_2(p_t,\pi)$ against the contraction bound.}
\label{tab:contraction}
\end{table}

\paragraph{How tight is the bound?} At $t=0$ the bound equals the true distance by construction. After that the
estimate falls faster than the bound, and the gap widens: the bound is $1.6$ times the estimate at $t=0.9$ and
over $5$ times at $t=2.7$. So the bound has the right form (exponential decay) but is loose for this $p_0$. The
reason is that the coupling argument in 2(c) only uses how fast two particles approach each other, which is
rate $\theta$. For our symmetric $p_0$ the mean is already $0$ (1(a)), so the actual distance to $\pi$ shrinks faster than
this worst case.

After $t\approx3$ the ratio falls again and drops below $1$ near $T$, i.e.\ the estimate appears to break the bound.
This is not a real violation. The true $W_2(p_t,\pi)$ keeps decreasing, but the estimate cannot go below its noise
floor of about $0.02$ (see 2(a)), while the bound continues down to $0.01$. Checking the bound at such small distances
would require many more samples.

\subsubsection*{2(e) Collapse against $\theta t$}

We repeat 2(a) for $\theta\in\{0.5,1,2\}$, sampling each run at the same values of $\theta t$ (so at $t = \theta t/\theta$).
Figure~\ref{fig:w2collapse} plots all six curves against $\theta t$. For each distance, the three values of $\theta$
fall on the same curve: at any $\theta t$ they differ by at most $0.016$ for $W_2(p_t,\pi)$ and $0.013$ for
$W_2(p_t,p_0)$, which is the size of the sampling noise. This confirms once more that the law of $X_t$ depends on
$\theta$ and $t$ only through $\theta t$ (1(d)), so every distance computed from it does too.

\begin{figure}[H]
\centering
\includegraphics[width=0.8\linewidth]{e_collapse.png}
\caption{$W_2(p_t,\pi)$ (blue) and $W_2(p_t,p_0)$ (orange) for $\theta\in\{0.5,1,2\}$ against rescaled time $\theta t$.}
\label{fig:w2collapse}
\end{figure}

\section*{Part 3: Score Function and the Stein Identity}

The score of $p_t$ is $s_t(x) = \partial_x\log p_t(x)$.

\subsection*{Derivations}

\subsubsection*{3(a) Fokker--Planck equation}

For an SDE $\dd X_t = f(X_t,t)\,\dd t + g(t)\,\dd W_t$, the density $p_t$ of $X_t$ evolves by the Fokker--Planck equation
\[
\partial_t p_t(x) = -\partial_x\bigl(f(x,t)\,p_t(x)\bigr) + \tfrac12 g(t)^2\,\partial_x^2 p_t(x).
\]
For~\eqref{eq:ou}, $f(x,t) = -\theta x$ and $g(t)=\sqrt{2\theta}$, so $\tfrac12g^2=\theta$ and
\[
\partial_t p_t(x) = -\partial_x\bigl(-\theta x\,p_t(x)\bigr) + \theta\,\partial_x^2 p_t(x)
= \theta\bigl(\partial_x(x\,p_t(x)) + \partial_x^2 p_t(x)\bigr).
\]
Bringing the right-hand side to the left,
\[
\partial_t p_t + \partial_x\Bigl(-\theta\bigl(x\,p_t + \partial_x p_t\bigr)\Bigr) = 0 ,
\]
which is the flux form $\partial_t p_t + \partial_x J_t = 0$ with
\[
\boxed{J_t(x) = -\theta\bigl(x\,p_t(x) + \partial_x p_t(x)\bigr).}
\]
\paragraph{Stationary flux.} For $\pi(x) = \frac{1}{\sqrt{2\pi}}e^{-x^2/2}$,
\[
\partial_x\pi(x) = \frac{1}{\sqrt{2\pi}}\bigl(-x\,e^{-x^2/2}\bigr) = -x\,\pi(x),
\]
so
\[
J_\pi(x) = -\theta\bigl(x\,\pi(x) + \partial_x\pi(x)\bigr) = -\theta\bigl(x\,\pi(x) - x\,\pi(x)\bigr) = 0
\]
for every $x$. The flux vanishes identically, so $\partial_t\pi = -\partial_x J_\pi = 0$ and $\pi$ is stationary, in agreement with 1(a).
Note that $J_t = -\theta\,p_t\,(x + s_t(x))$, so $J_t\equiv0$ exactly when $s_t(x) = -x$, which is the score of $\pi$.

\subsubsection*{3(b) Score in closed form}

From 1(c), $p_t(x) = \tfrac12\varphi_1(x) + \tfrac12\varphi_2(x)$ with
\[
\varphi_1(x) = \varphi(x\mid -m_t, v_t),\qquad \varphi_2(x) = \varphi(x\mid m_t, v_t),\qquad m_t = 3e^{-\theta t},\quad v_t = 1-0.75e^{-2\theta t}.
\]
By the chain rule, $s_t(x) = \partial_x\log p_t(x) = \partial_x p_t(x)/p_t(x)$, and
$\partial_x p_t = \tfrac12\partial_x\varphi_1 + \tfrac12\partial_x\varphi_2$.
For a Gaussian density $\varphi(x\mid m,v) = \frac{1}{\sqrt{2\pi v}}\exp\!\bigl(-\frac{(x-m)^2}{2v}\bigr)$,
\[
\partial_x\varphi(x\mid m,v) = \frac{1}{\sqrt{2\pi v}}\Bigl(-\frac{x-m}{v}\Bigr)\exp\!\Bigl(-\frac{(x-m)^2}{2v}\Bigr) = -\frac{x-m}{v}\,\varphi(x\mid m,v).
\]
Hence $\partial_x\varphi_k = s_{k,t}\,\varphi_k$ with component scores
\[
s_{1,t}(x) = -\frac{x+m_t}{v_t},\qquad s_{2,t}(x) = -\frac{x-m_t}{v_t}.
\]
Substituting,
\[
s_t(x) = \frac{\tfrac12 s_{1,t}\varphi_1 + \tfrac12 s_{2,t}\varphi_2}{\tfrac12\varphi_1 + \tfrac12\varphi_2}
= w_{1,t}(x)\,s_{1,t}(x) + w_{2,t}(x)\,s_{2,t}(x),
\]
with
\[
w_{1,t}(x) = \frac{\varphi_1(x)}{\varphi_1(x)+\varphi_2(x)},\qquad w_{2,t}(x) = \frac{\varphi_2(x)}{\varphi_1(x)+\varphi_2(x)} .
\]
The weights are non-negative and sum to one, so $s_t$ is a weighted average of the two component scores.
By Bayes' rule, with the label $Z$ from 1(c),
\[
w_{k,t}(x) = \frac{P(Z=k)\,p(x\mid Z=k)}{p_t(x)} = P(Z=k\mid X_t = x),
\]
the posterior probability that a particle found at $x$ came from component $k$. Therefore
\[
\boxed{s_t(x) = w_{1,t}(x)\,s_{1,t}(x) + w_{2,t}(x)\,s_{2,t}(x).}
\]

As $t\to\infty$, $m_t\to0$ and $v_t\to1$, so both component scores tend to $-x$, and hence so does $s_t(x)$: the
score of $\pi$.

\subsubsection*{3(c) The score stays finite}

Taken literally, $\sup_x|s_t(x)|$ is infinite for every $t$: the score grows linearly in $|x|$, and even the limiting
score $-x$ is unbounded. What does stay finite is the score at every point and on every bounded set of $x$,
uniformly over all $t\ge0$, including $t=0$. This is the property we prove.

By 3(b), $s_t(x)$ is a weighted average of $s_{1,t}(x)$ and $s_{2,t}(x)$ with weights in $[0,1]$ that sum to one. A weighted
average is never larger in size than the larger of the two values being averaged, so
\[
|s_t(x)| \le \max\bigl(|s_{1,t}(x)|,\,|s_{2,t}(x)|\bigr) = \max\Bigl(\frac{|x+m_t|}{v_t},\,\frac{|x-m_t|}{v_t}\Bigr) \le \frac{|x| + m_t}{v_t}.
\]
For all $t\ge0$ we have $m_t\le3$ and $v_t = 1 - 0.75e^{-2\theta t}\ge 0.25$, so
\[
|s_t(x)| \le 4|x| + 12 \qquad\text{for all } x\in\mathbb R,\ t\ge0 .
\]
Hence $\sup_{|x|\le L}|s_t(x)| \le 4L + 12$ for every $L$, with a bound that does not depend on $t$.

\paragraph{Which feature of $p_0$ guarantees this.} The key quantity is the lower bound $v_t\ge 0.25$, which comes
from the components of $p_0$ having positive variance $0.5^2$. The variance of each component starts at $0.25$
and only grows, so the $1/v_t$ factors in the score never blow up.

\paragraph{Two point masses.} If $p_0 = \tfrac12\delta_{-3}+\tfrac12\delta_{3}$, the same computation as in 1(c) gives
components $\N(\mp3e^{-\theta t}, v_t)$ with $v_t = 1-e^{-2\theta t}$, so $v_t\to0$ as $t\to0$.
\begin{itemize}
\item At $t=0$, $p_0$ has no density, so $\log p_0$ and the score are not defined at all.
\item For small $t>0$, $v_t\approx 2\theta t$, so both component scores $-(x\pm m_t)/v_t$ are huge. Take for example $x=1$.
It is much closer to the right component, so $w_{2,t}(1)\to1$ and $s_t(1)\approx s_{2,t}(1) = (m_t-1)/v_t\approx 2/(2\theta t)\to\infty$.
The same happens at every $x\ne\pm3$.
\end{itemize}
So with point masses the score is finite for each $t>0$ but not uniformly in $t$: it blows up like $1/t$ as $t\to0$. This
matters for diffusion models, where the score must be learned and evaluated down to small $t$.

\subsection*{Implementation}

\subsubsection*{Density and score over time}

Figure~\ref{fig:score} plots $p_t$ and the closed-form score $s_t$ from 3(b) on a common grid at $t\in\{0,0.5,1,2,T\}$.

\begin{figure}[H]
\centering
\includegraphics[width=\linewidth]{score_plots.png}
\caption{Density $p_t$ (left) and score $s_t$ (right) at several times; the dashed line is $-x$, the score of $\pi$.}
\label{fig:score}
\end{figure}

The score points uphill on the density: it is positive where $p_t$ increases and negative where it decreases,
and zero at every peak and trough.
\begin{itemize}
\item \textbf{$t=0$.} Inside each bump the score is close to that bump's own score, a steep line
$-(x\pm3)/0.25$ through the bump's centre. Near $x=0$ the weights in 3(b) switch from one component to the other,
so the score jumps sharply from about $-11$ to $+11$. This S-shaped jump marks the trough between the two modes.
\item \textbf{$t=0.5$ and $t=1$.} The components move inwards and widen, so the lines become less steep and the
jump at the origin becomes smaller and smoother. While the density is bimodal, the score still crosses zero three
times (at the two peaks and at the trough).
\item \textbf{After $t^*\approx1.14$.} The density has a single peak, so the score crosses zero only once, at $x=0$.
By $t=2$ it is already close to a straight line.
\item \textbf{Limit.} As $t\to\infty$, $m_t\to0$ and $v_t\to1$, so the score tends to
\[
s_\infty(x) = -x,
\]
the score of $\pi=\N(0,1)$. At $t=T$ it is indistinguishable from $-x$ on the plot.
\end{itemize}

\subsubsection*{Stein identity}

\paragraph{Statement and proof.} Let $f$ be differentiable with
\[
f(x)\,p_t(x)\to0 \text{ as } |x|\to\infty, \qquad \E_{p_t}|f'(X)|<\infty, \qquad \E_{p_t}|f(X)\,s_t(X)|<\infty .
\]
Since $s_t = \partial_x p_t/p_t$, we have $f\,s_t\,p_t = f\,\partial_x p_t$. Integrating by parts,
\[
\E_{p_t}[f'(X)] = \int_{-\infty}^{\infty} f'(x)\,p_t(x)\,\dd x
= \Bigl[f(x)\,p_t(x)\Bigr]_{-\infty}^{\infty} - \int_{-\infty}^{\infty} f(x)\,\partial_x p_t(x)\,\dd x
= -\E_{p_t}[f(X)\,s_t(X)] ,
\]
where the boundary term vanishes by the decay condition. Because $p_t$ has Gaussian tails, the condition holds for
any $f$ that grows at most polynomially, which includes all four test functions.

\paragraph{Numerical check.} We use the $N=2\times10^4$ particles of the forward simulation at
$t\in\{0, 1.35, 2.70, 4.05, 5.40\}$ and evaluate the exact score $s_t$ from 3(b) at each particle. For each test function we
report the Stein residual
\[
R_f = \Bigl|\frac1N\sum_{i=1}^N f'(X_i) + \frac1N\sum_{i=1}^N f(X_i)\,s_t(X_i)\Bigr|
\]
and its standard error SE, the sample standard deviation of $f'(X_i) + f(X_i)s_t(X_i)$ divided by $\sqrt N$.
A check passes when $R_f \le 2\,\mathrm{SE}$.

\begin{table}[H]
\centering
\small
\begin{tabular}{c|cc|cc|cc|cc}
\hline
& \multicolumn{2}{c|}{$f=x$} & \multicolumn{2}{c|}{$f=x^2$} & \multicolumn{2}{c|}{$f=\sin x$} & \multicolumn{2}{c}{$f=e^{-x^2/2}$} \\
$t$ & $R_f$ & SE & $R_f$ & SE & $R_f$ & SE & $R_f$ & SE \\
\hline
0.00 & 0.0565 & 0.0431 & 0.0712 & 0.1422 & 0.0078 & 0.0090 & 0.0001 & 0.0010 \\
1.35 & 0.0019 & 0.0111 & 0.0194 & 0.0259 & 0.0044 & 0.0067 & 0.0027 & 0.0047 \\
2.70 & 0.0078 & 0.0102 & 0.0041 & 0.0203 & 0.0006 & 0.0070 & 0.0052 & 0.0061 \\
4.05 & 0.0020 & 0.0100 & 0.0259 & 0.0186 & 0.0048 & 0.0071 & 0.0000 & 0.0062 \\
5.40 & 0.0057 & 0.0102 & 0.0244 & 0.0191 & 0.0029 & 0.0071 & 0.0018 & 0.0062 \\
\hline
\end{tabular}
\caption{Stein residuals $R_f$ and standard errors. Every $R_f$ is below $2\,\mathrm{SE}$; the largest ratio is $R_f/\mathrm{SE} = 1.39$.}
\label{tab:stein}
\end{table}

All $20$ checks pass. The ratio $R_f/\mathrm{SE}$ ranges from $0.00$ to $1.39$, which is what we expect for a quantity
whose true mean is zero. The standard errors are largest at $t=0$, where the score is steep ($|s_0|$ reaches about $11$),
so the terms $f(X_i)s_t(X_i)$ vary the most.

\section*{Part 4: Reverse Diffusion with the Exact Score}

For a forward SDE $\dd X_t = f(X_t,t)\,\dd t + g(t)\,\dd W_t$ with marginals $p_t$, Anderson's time reversal states that the
reverse-time process satisfies
\begin{equation}\label{eq:anderson}
\dd X_t = \bigl[f(X_t,t) - g(t)^2 s_t(X_t)\bigr]\dd t + g(t)\,\dd\bar W_t ,
\end{equation}
where $\bar W$ is a Brownian motion running backwards in time and $\dd t<0$.

\subsection*{Derivations}

\subsubsection*{4(a) Reverse SDE for the OU process}

For~\eqref{eq:ou}, $f(x,t) = -\theta x$ and $g(t) = \sqrt{2\theta}$, so $g^2 = 2\theta$. Substituting into~\eqref{eq:anderson},
\[
\dd X_t = \bigl[-\theta X_t - 2\theta\,s_t(X_t)\bigr]\dd t + \sqrt{2\theta}\,\dd\bar W_t
= -\theta\bigl[X_t + 2s_t(X_t)\bigr]\dd t + \sqrt{2\theta}\,\dd\bar W_t .
\]
The reverse drift is
\[
\boxed{f_{\mathrm{rev}}(x,t) = -\theta\bigl[x + 2s_t(x)\bigr].}
\]

\subsubsection*{4(b) Discrete reverse update}

We step from time $t$ down to $t-\Delta t$ with $\Delta t>0$, so
\[
\dd X_t \approx X_{t-\Delta t} - X_t, \qquad \dd t = -\Delta t, \qquad \dd\bar W_t \approx \bar W_{t-\Delta t} - \bar W_t \sim \N(0,\Delta t).
\]
The Brownian increment has variance equal to the length of the time step, so we write $\dd\bar W_t = \sqrt{\Delta t}\,Z$ with
$Z\sim\N(0,1)$; its sign does not matter because $Z$ is symmetric. Substituting into 4(a),
\[
X_{t-\Delta t} - X_t = -\theta\bigl[X_t + 2s_t(X_t)\bigr](-\Delta t) + \sqrt{2\theta\Delta t}\,Z ,
\]
which gives
\begin{equation}\label{eq:reverse}
\boxed{X_{t-\Delta t} = X_t + \bigl[\theta X_t + 2\theta\,s_t(X_t)\bigr]\Delta t + \sqrt{2\theta\Delta t}\,Z .}
\end{equation}
The negative $\dd t$ flips the sign of \emph{both} drift terms.

\paragraph{Note on the sign in the assignment.} The assignment's equation (9) has $-2\theta s_t(X_t)$ in place of
$+2\theta s_t(X_t)$. That version is incorrect. To see this, take $p_t=\pi$, so $s_t(x)=-x$; the reverse process must then
leave $\pi$ unchanged.
\begin{itemize}
\item With~\eqref{eq:reverse}: $X_{t-\Delta t} = (1-\theta\Delta t)X_t + \sqrt{2\theta\Delta t}\,Z$, whose variance is
$(1-\theta\Delta t)^2 + 2\theta\Delta t = 1 + \theta^2\Delta t^2$. So $\pi$ is preserved up to $O(\Delta t^2)$.
\item With the printed minus sign: $X_{t-\Delta t} = (1+3\theta\Delta t)X_t + \sqrt{2\theta\Delta t}\,Z$. Every step multiplies by
more than $1$, and the samples blow up.
\end{itemize}
All results below use~\eqref{eq:reverse}.

\subsubsection*{4(c) Stability near a mode}

Near one mode of $p_0$ the score is approximately that of a single Gaussian, $s_t(x)\approx-(x-m)/v$, where $m$ is the
centre and $v$ the variance of the mode. Write the deviation from the mode as $y_t = X_t - m$, so $X_t = y_t + m$.
Substituting into~\eqref{eq:reverse},
\[
y_{t-\Delta t} + m = (y_t+m) + \Bigl[\theta(y_t+m) - 2\theta\,\frac{y_t}{v}\Bigr]\Delta t + \sqrt{2\theta\Delta t}\,Z ,
\]
so
\[
y_{t-\Delta t} = \underbrace{\Bigl[1 - \theta\Bigl(\frac2v - 1\Bigr)\Delta t\Bigr]}_{A}\,y_t + \theta m\,\Delta t + \sqrt{2\theta\Delta t}\,Z .
\]
The term $\theta m\Delta t$ is a constant shift and the noise term does not depend on $y_t$, so neither affects whether
deviations grow. The recursion does not amplify $y_t$ when $|A|\le1$, i.e.\ $-1\le A\le1$.
\begin{itemize}
\item \textbf{Upper inequality} $A\le1$: $\theta\bigl(\frac2v-1\bigr)\Delta t\ge0$, which holds for every $\Delta t>0$ when $v<2$.
\item \textbf{Lower inequality} $A\ge-1$: $\theta\bigl(\frac2v-1\bigr)\Delta t\le2$, i.e.
\[
\boxed{\Delta t \le \frac{2}{\theta\bigl(\frac2v-1\bigr)}.}
\]
\end{itemize}
At $v=v_0=0.25$ and $\theta=1$, $\frac2v = 8$ and
\[
\Delta t \le \frac{2}{8-1} = \frac27 \approx 0.286 .
\]
The narrowest modes (smallest $v$) occur at $t=0$, so this is the most restrictive value. The default $\Delta t = 10^{-3}$
is far below it.

\subsection*{Implementation}

\subsubsection*{4(a) and 4(c) Running the reverse SDE}

We implement~\eqref{eq:reverse} with the exact score from 3(b), stepping $t = T, T-\Delta t, \dots, \Delta t$ with
$\Delta t = 10^{-3}$ and $T = 5.40$, and evaluating the score at the current time $t$ in each step. The $N=2\times10^4$ particles at
the end form the empirical law $\hat p_0$.

\subsubsection*{4(b) Starting from $\pi$ instead of $p_T$}

We start the reverse process from $N$ samples of $\pi=\N(0,1)$ rather than from $p_T$. This is justified by Part 2:
the contraction bound gives $W_2(p_T,\pi)\le e^{-\theta T}W_2(p_0,\pi)$, and $T$ was chosen in 2(d) so that this is at most
$10^{-2}$. So the starting distribution differs from the correct one, $p_T$, by at most $0.01$ in $W_2$. This is the
error the substitution introduces.

\subsubsection*{4(d) Comparing $\hat p_0$ with $p_0$}

\begin{figure}[H]
\centering
\includegraphics[width=0.8\linewidth]{reverse_sde_p0.png}
\caption{Histogram of $\hat p_0$ from the reverse SDE ($\Delta t=10^{-3}$) against the true $p_0$.}
\label{fig:reverse}
\end{figure}

Figure~\ref{fig:reverse} shows that the reverse SDE recovers both modes, with the right locations, heights and widths.
Starting from a single Gaussian, the reverse dynamics separates the particles into the two original modes.

\begin{table}[H]
\centering
\begin{tabular}{lcc}
\hline
& $\hat p_0$ & $p_0$ \\
\hline
mean & $0.005$ & $0$ \\
variance & $9.255$ & $9.25$ \\
\hline
\end{tabular}
\caption{First two moments. The true variance is $3^2+0.5^2=9.25$.}
\label{tab:reverse}
\end{table}

Both moments agree with $p_0$ to within sampling error (the standard errors are about $0.02$ for the mean and $0.02$ for
the variance). The estimated distance is $W_2(\hat p_0,p_0) = 0.188$. For comparison, two independent samples of size $N$
drawn directly from $p_0$ are $0.227$ apart under the same estimator, so $\hat p_0$ is as close to $p_0$ as a sample
from $p_0$ itself. These values are large because of how the estimator handles the two modes: the number of points
in each mode varies randomly by about $\sqrt N/2\approx70$ between samples, and after sorting, these surplus points
are paired with points in the other mode, about $6$ away.

\subsubsection*{4(e) Convergence in $\Delta t$}

\begin{table}[H]
\centering
\begin{tabular}{cc}
\hline
$\Delta t$ & $W_2(\hat p_0,p_0)$ \\
\hline
$10^{-3}$ & $0.188$ \\
$5\times10^{-4}$ & $0.315$ \\
$2.5\times10^{-4}$ & $0.149$ \\
\hline
\end{tabular}
\caption{$W_2(\hat p_0,p_0)$ for the reverse SDE as $\Delta t$ is halved and quartered.}
\label{tab:stepsize}
\end{table}

The values do not decrease with $\Delta t$: the ratios give observed orders $\log_2(0.188/0.315) = -0.74$ and
$\log_2(0.315/0.149) = 1.08$, with no consistent pattern. So no order of convergence can be identified at these step sizes.
The reason is that all three values are of the same size as the $W_2$ distance between two independent samples of
$p_0$ ($0.227$, see 4(d)). The discretisation error of Euler--Maruyama at $\Delta t\le10^{-3}$ is much smaller than this sampling
noise, so it cannot be seen. The theoretical expectation is weak order $1$, i.e.\ an error proportional to $\Delta t$.

\subsubsection*{4(f) Probability flow ODE}

The probability flow ODE $\dd X_t/\dd t = -\theta X_t - \theta s_t(X_t)$ (equation (10) of the assignment) is discretised with the explicit Euler method, again with $\dd t=-\Delta t$:
\[
X_{t-\Delta t} = X_t + \theta\bigl[X_t + s_t(X_t)\bigr]\Delta t .
\]
It starts from the same $N$ samples of $\pi$ as the SDE and has no noise term, so each particle follows a deterministic
path. Figure~\ref{fig:flowode} compares $W_2(\hat p_0,p_0)$ for the two schemes at equal $\Delta t$, from $10^{-3}$ up to the coarse
step $0.2$ (still below the stability limit $2/7$ from 4(c)).

\begin{figure}[H]
\centering
\includegraphics[width=0.8\linewidth]{flow_ode_vs_sde.png}
\caption{$W_2(\hat p_0,p_0)$ against $\Delta t$ for the reverse SDE and the probability flow ODE.}
\label{fig:flowode}
\end{figure}

\begin{itemize}
\item \textbf{ODE.} The curve is smooth: since there is no noise and every run starts from the same particles, the only
thing that changes with $\Delta t$ is the discretisation error. For $\Delta t\le0.05$ it is flat at about $0.26$, the sampling floor of
this particular set of particles. At $\Delta t = 0.1$ and $0.2$ it rises to $0.27$ and $0.35$, so the discretisation error becomes
visible at coarse steps.
\item \textbf{SDE.} Each run uses new noise, so the number of particles that end up in each mode changes from run to run,
and the curve jumps between $0.09$ and $0.35$ with no trend in $\Delta t$. Even at $\Delta t=0.2$, its value ($0.28$) lies within
this run-to-run range.
\end{itemize}

At the coarsest step, $\Delta t=0.2$, the SDE gives the smaller distance ($0.28$ against $0.35$). But since the SDE values
vary by about $\pm0.1$ between runs, a single run per step size cannot show reliably which scheme is more accurate at coarse
steps. What the plot does show is that the ODE's error grows clearly once $\Delta t$ exceeds about $0.05$, while for
$\Delta t\le0.05$ both schemes are limited by sampling noise rather than by the step size.

\section*{Part 5: Effect of Imperfect Score Estimation}

We repeat the Part 4 reconstruction (same $N$, $\Delta t$, $T$, seed, start from $\pi$) with an estimated score $\hat s_t$ in~\eqref{eq:reverse}:
\[
X_{t-\Delta t} = X_t + \bigl[\theta X_t + 2\theta\,\hat s_t(X_t)\bigr]\Delta t + \sqrt{2\theta\Delta t}\,Z .
\]
For each run we record $W_2(\hat p_0,p_0)$ and the Fisher divergence
$D_F = \frac1T\int_0^T \E_{p_t}\bigl[(\hat s_t(X)-s_t(X))^2\bigr]\dd t$, where $\E_{p_t}$ is the average over the Part 1 forward particles at
55 equally spaced times and the time integral uses the trapezoidal rule. As in 4(d), $W_2\approx0.2$ is the noise floor of the
estimator.

\subsection*{Case A: Random error}

On a grid of $K=200$ points in $[-6,6]$ we draw $e_k\sim\N(0,\sigma_\epsilon^2)$ once, and set $\hat s_t = s_t + e$, with $e$ the linear interpolation of
the $e_k$, held fixed for the whole run.

\begin{table}[H]
\centering
\begin{tabular}{ccccc}
\hline
$\sigma_\epsilon$ & $W_2(\hat p_0,p_0)$ & $D_F$ & mean & variance \\
\hline
0.1 & 0.315 & 0.005 & $-0.030$ & 9.244 \\
0.3 & 0.445 & 0.047 & $-0.084$ & 9.212 \\
1.0 & 0.875 & 0.527 & $-0.304$ & 9.054 \\
3.0 & 1.649 & 4.740 & $-0.872$ & 8.316 \\
\hline
\end{tabular}
\caption{Case A. The true mean and variance are $0$ and $9.25$.}
\label{tab:caseA}
\end{table}

$W_2$ grows with $\sigma_\epsilon$, and $D_F$ grows like $\sigma_\epsilon^2$. A single fixed draw of $e$ is not zero on average where the
particles are, so it acts locally like a bias; in this draw it shifts the mean to the left.

\subsection*{Case B: Systematic bias}

Here $\hat s_t = s_t + c$, so $D_F = c^2$ exactly.

\paragraph{Prediction.} Take a single Gaussian mode, $p_t = \N(m_t,v_t)$ with $m_t = m_0e^{-\theta t}$ and $v_t = 1-(1-v_0)e^{-2\theta t}$, so
$s_t(x) = -(x-m_t)/v_t$.
\begin{enumerate}
\item Substituting $\hat s_t = s_t + c$ into the update and taking expectations, with $\mu_t = \E[X_t]$,
\[
\mu_{t-\Delta t} = \mu_t + \theta\mu_t\Delta t + 2\theta\,\E[s_t(X_t)]\,\Delta t + 2\theta c\,\Delta t,
\qquad \E[s_t(X_t)] = -\frac{\mu_t - m_t}{v_t}.
\]
This expectation is $0$ only when $\mu_t = m_t$, i.e.\ when the particles follow $p_t$. Once the bias has moved them, the true score
pulls them back, so the term must be kept.
\item The correct mean satisfies the same recursion with $c=0$ and $\mu_t = m_t$: $\ m_{t-\Delta t} = m_t + \theta m_t\Delta t$.
\item Subtracting, the shift $\delta_t = \mu_t - m_t$ satisfies
\[
\delta_{t-\Delta t} = \delta_t + \theta\Bigl(1-\frac2{v_t}\Bigr)\delta_t\Delta t + 2\theta c\,\Delta t .
\]
\item Rearranging, $(\delta_t - \delta_{t-\Delta t})/\Delta t = \theta(2/v_t - 1)\delta_t - 2\theta c$. Letting $\Delta t\to0$,
\[
\frac{\dd\delta}{\dd t} = \theta\Bigl(\frac2{v_t}-1\Bigr)\delta - 2\theta c, \qquad \delta_T = 0 ,
\]
since the run starts from $\pi\approx p_T$ at $t=T$.
\item Integrating factor $e^{-\theta t}/v_t$. Using $\dd v_t/\dd t = 2\theta(1-v_t)$,
\[
\frac{\dd}{\dd t}\Bigl(\frac{e^{-\theta t}}{v_t}\Bigr) = \frac{e^{-\theta t}}{v_t}\Bigl(-\theta - \frac{2\theta(1-v_t)}{v_t}\Bigr)
= -\theta\Bigl(\frac2{v_t}-1\Bigr)\frac{e^{-\theta t}}{v_t},
\]
so
\[
\frac{\dd}{\dd t}\Bigl(\frac{e^{-\theta t}}{v_t}\,\delta_t\Bigr) = \frac{e^{-\theta t}}{v_t}\Bigl(\frac{\dd\delta}{\dd t} - \theta\Bigl(\frac2{v_t}-1\Bigr)\delta_t\Bigr)
= -2\theta c\,\frac{e^{-\theta t}}{v_t} .
\]
\item Integrating from $0$ to $T$ and using $\delta_T = 0$: $\ -\delta_0/v_0 = -2\theta c\int_0^T e^{-\theta s}/v_s\,\dd s$, i.e.
\[
\boxed{\delta_0 = 2\theta c\,v_0\int_0^T \frac{e^{-\theta s}}{v_s}\,\dd s \approx 0.76\,c}
\qquad (v_0 = 0.25,\ \theta=1,\ T=5.40).
\]
\end{enumerate}
A constant bias shifts the mean of the reconstruction in the direction of $c$, by an amount proportional to $c$.

\begin{table}[H]
\centering
\begin{tabular}{ccccc}
\hline
$c$ & $W_2(\hat p_0,p_0)$ & $D_F$ & mean & variance \\
\hline
$-1.0$ & 4.540 & 1.00 & $-3.692$ & 0.534 \\
$-0.5$ & 3.637 & 0.25 & $-2.673$ & 3.901 \\
$-0.2$ & 2.235 & 0.04 & $-1.252$ & 8.013 \\
$0.2$ & 2.186 & 0.04 & $1.265$ & 7.974 \\
$0.5$ & 3.602 & 0.25 & $2.692$ & 3.806 \\
$1.0$ & 4.497 & 1.00 & $3.700$ & 0.499 \\
\hline
\end{tabular}
\caption{Case B.}
\label{tab:caseB}
\end{table}

\paragraph{Test.} The direction is confirmed: the mean moves with the sign of $c$, symmetrically. The size is much larger
than $0.76\,c$ (e.g.\ $1.26$ at $c=0.2$) and not linear in $c$. The prediction covers a single mode; with two modes the bias also
moves particles from one mode to the other. A mean of $1.26$ corresponds to about $71\%$ of the particles in the right mode,
and at $|c|=1$ almost all particles end up in one mode, which is why the variance drops to about $0.5$.

\subsection*{Case C: Multiplicative error}

Here $\hat s_t = (1+\alpha)s_t$, so $D_F = \alpha^2\,\frac1T\int_0^T\E_{p_t}[s_t^2]\,\dd t\approx\alpha^2$.

\paragraph{Prediction.} Take the same single Gaussian mode, so $\hat s_t(x) = -(1+\alpha)(x-m_t)/v_t$.
\begin{enumerate}
\item Let $y_t = X_t - m_t$. Substituting into the update and using $m_{t-\Delta t} = m_t + \theta m_t\Delta t$ (step 2 of Case B), the
$m_t$ terms cancel:
\[
y_{t-\Delta t} = (1 + a_t\Delta t)\,y_t + \sqrt{2\theta\Delta t}\,Z, \qquad a_t = \theta\Bigl(1-\frac{2(1+\alpha)}{v_t}\Bigr).
\]
\item The mean of $y$ starts at $0$ and stays $0$. For the variance $w_t = \E[y_t^2]$, using that $Z$ is independent of $y_t$ with $\E Z^2=1$,
\[
w_{t-\Delta t} = (1+a_t\Delta t)^2 w_t + 2\theta\Delta t = w_t + 2a_tw_t\Delta t + 2\theta\Delta t + O(\Delta t^2).
\]
\item Letting $\Delta t\to0$,
\[
\frac{\dd w}{\dd t} = -2a_tw - 2\theta = 2\theta\Bigl(\frac{2(1+\alpha)}{v_t}-1\Bigr)w - 2\theta, \qquad w_T = 1 .
\]
\item Try $w_t = v_t/(1+2\alpha)$. With $\dd v_t/\dd t = 2\theta(1-v_t)$:
\[
\text{left} = \frac{2\theta(1-v_t)}{1+2\alpha}, \qquad
\text{right} = \frac{2\theta\bigl(2(1+\alpha)-v_t\bigr) - 2\theta(1+2\alpha)}{1+2\alpha} = \frac{2\theta(1-v_t)}{1+2\alpha} .
\]
So it is a solution.
\item The difference $u_t = w_t - v_t/(1+2\alpha)$ solves the equation without the $-2\theta$ term, so
$u_0 = u_T\exp\bigl(-\int_0^T 2\theta\,(2(1+\alpha)/v_s-1)\,\dd s\bigr)$. Since $v_s\le1$, the exponent is at most $-2\theta(1+2\alpha)T$, so for $\alpha>-\tfrac12$
the difference is negligible after a long run:
\[
\boxed{\Var(\hat p_0) \approx \frac{v_0}{1+2\alpha}, \qquad \alpha>-\tfrac12 .}
\]
\end{enumerate}
A score that is too strong ($\alpha>0$) shrinks the variance; one that is too weak ($\alpha<0$) inflates it, and for $\alpha\le-\tfrac12$ the
variance is not controlled.

\begin{table}[H]
\centering
\begin{tabular}{cccccc}
\hline
$\alpha$ & $W_2(\hat p_0,p_0)$ & $D_F$ & mean & variance & predicted $9+0.25/(1+2\alpha)$ \\
\hline
$-0.50$ & 1.344 & 0.250 & 0.028 & 18.004 & -- \\
$-0.25$ & 0.216 & 0.063 & 0.021 & 10.090 & 9.50 \\
$-0.10$ & 0.177 & 0.010 & 0.008 & 9.442 & 9.31 \\
$0.10$ & 0.203 & 0.010 & 0.004 & 9.162 & 9.21 \\
$0.25$ & 0.260 & 0.063 & $-0.003$ & 9.091 & 9.17 \\
$0.50$ & 0.257 & 0.250 & 0.003 & 9.062 & 9.13 \\
\hline
\end{tabular}
\caption{Case C. For two modes at $\pm3$, the total variance is about $9$ plus the within-mode variance $0.25/(1+2\alpha)$.}
\label{tab:caseC}
\end{table}

\paragraph{Test.} The mean stays at $0$. The variance moves in the predicted direction for every $\alpha$: above $9.25$ for $\alpha<0$ and
below for $\alpha>0$. The size is close for $\alpha>0$ and underestimated for $\alpha<0$, where a weaker score also separates the two
modes less cleanly; at $\alpha=-0.5$ the variance reaches $18.0$.

\subsection*{Comparing the cases}

\begin{figure}[H]
\centering
\includegraphics[width=0.8\linewidth]{w2_vs_fisher.png}
\caption{$W_2(\hat p_0,p_0)$ against $D_F$ for all runs.}
\label{fig:fisher}
\end{figure}

\begin{figure}[H]
\centering
\includegraphics[width=\linewidth]{w2_vs_parameter.png}
\caption{$W_2(\hat p_0,p_0)$ against the native error parameter of each case.}
\label{fig:native}
\end{figure}

$D_F$ alone does not determine the reconstruction error. At similar $D_F$, the bias is by far the most harmful
($W_2\approx2.2$ at $D_F=0.04$), against $0.45$ for the random error ($D_F=0.047$) and about $0.2$ for the multiplicative error ($D_F=0.063$).
A bias pushes every particle the same way and moves probability between the modes, $6$ apart. A multiplicative error
keeps the drift symmetric, so only the spread of each mode changes. The random error lies in between.

\section*{Part 6: Application Challenge, Noisy Sensor Measurements}

We observe $Y^{(i)} = X_0^{(i)} + \eta^{(i)}$, $i=1,\dots,N$, with $X_0^{(i)}\sim p_0$ and $\eta^{(i)}\sim\N(0,\sigma^2)$ independent, where $\sigma=\sigma_{\mathrm{sensor}}$.
The target is $p_0$.

\subsection*{Derivations}

\subsubsection*{6(a) Rescaling to a forward marginal}

By 1(b), $X_t = e^{-\theta t}X_0 + \sqrt{1-e^{-2\theta t}}\,Z$ with $Z\sim\N(0,1)$ independent of $X_0$. For $a>0$,
\[
aY = aX_0 + a\eta, \qquad a\eta\sim\N(0,a^2\sigma^2) \text{ independent of } X_0 .
\]
Both are ``$X_0$ scaled, plus independent Gaussian noise'', so $aY$ has the law $p_{t^*}$ (for any $p_0$) if the scale and the noise
variance match:
\[
a = e^{-\theta t^*}, \qquad a^2\sigma^2 = 1 - e^{-2\theta t^*} .
\]
Substituting the first into the second, $e^{-2\theta t^*}\sigma^2 = 1 - e^{-2\theta t^*}$, i.e.\ $e^{-2\theta t^*}(1+\sigma^2) = 1$. Hence
\[
\boxed{t^* = \frac{\ln(1+\sigma^2)}{2\theta}, \qquad a = e^{-\theta t^*} = \frac{1}{\sqrt{1+\sigma^2}} .}
\]

\subsubsection*{6(b) Sign and behaviour of $t^*$}

For $\sigma>0$, $1+\sigma^2>1$, so $\ln(1+\sigma^2)>0$ and $t^*>0$. Since $\ln(1+\sigma^2)$ is increasing in $\sigma$, $t^*$ grows with the sensor noise: a
noisier sensor corresponds to a later, more diffused time. For small noise $t^*\approx\sigma^2/(2\theta)\to0$, and for large noise
$t^*\approx\ln\sigma/\theta\to\infty$, so $t^*$ grows only logarithmically. At the same time $a\to0$: the signal is scaled down more and more.

\subsubsection*{6(c) Why $Y$ cannot be used without the factor $a$}

The reverse process started at time $t$ is only correct if its starting points have law $p_t$. Without rescaling, $Y = X_0 + \eta$
keeps the signal $X_0$ at full scale, while every $X_t$ carries it scaled by $e^{-\theta t}$. For our $p_0$, $Y$ is a mixture of
$\N(\pm3,\,0.25+\sigma^2)$, whereas $p_t$ is a mixture of $\N(\pm3e^{-\theta t},\,v_t)$. Matching the mode positions requires $e^{-\theta t} = 1$, i.e.\ $t=0$,
but then $v_0 = 0.25 \ne 0.25 + \sigma^2$. So for $\sigma>0$ there is no time $t$ at which $Y$ has law $p_t$. Starting the reverse process
from $Y$ at any $t$ gives it inputs that are too wide (if $t$ is small) or too far apart (if $t$ is large), and the score of $p_t$ then
pushes them to the wrong places.

\subsubsection*{Task (2): Linear shrinkage $\lambda$}

Treat $X_0$ as Gaussian with the correct first two moments, $\E[X_0] = 0$ and $\Var(X_0) = s_0^2 = 3^2 + 0.5^2 = 9.25$. For
$\hat X = \lambda Y = \lambda X_0 + \lambda\eta$, using that $X_0$ and $\eta$ are independent with mean $0$,
\[
\E\bigl[(\lambda Y - X_0)^2\bigr] = \E\bigl[((\lambda-1)X_0 + \lambda\eta)^2\bigr] = (\lambda-1)^2s_0^2 + \lambda^2\sigma^2 .
\]
Setting the derivative in $\lambda$ to zero, $2(\lambda-1)s_0^2 + 2\lambda\sigma^2 = 0$, so
\[
\boxed{\lambda = \frac{s_0^2}{s_0^2+\sigma^2} = \frac{9.25}{9.25+\sigma^2} .}
\]
(Because the mean is $0$, no intercept is needed.)

\subsection*{Tasks}

\paragraph{Setup.} For each $N\in\{1000, 5000, 20000\}$ and $\sigma\in\{0.1, 0.3, 1, 3\}$ (a factor of $30$), we draw $N$ observations $Y$ and
compare three reconstructions:
\begin{enumerate}
\item the raw observations $Y$;
\item the shrinkage estimate $\lambda Y$;
\item reverse diffusion: $X_{t^*} = aY$, integrated with~\eqref{eq:reverse} and the exact score from $t^*$ down to $0$.
\end{enumerate}
\textbf{Assumptions.} For reverse diffusion we use steps of size $t^*/\lceil t^*/10^{-3}\rceil\le10^{-3}$, so that the run ends exactly at
$t=0$. For $\sigma=0.1$, $t^* = 0.005$, so this is only $5$ steps. Each reconstruction is compared by $W_2$ with a fixed reference sample of $N$
points from $p_0$, using one run per $(N,\sigma)$.

\begin{table}[H]
\centering
\begin{tabular}{cccccc}
\hline
$N$ & $\sigma$ & $t^*$ & raw $Y$ & shrinkage & reverse diffusion \\
\hline
1000 & 0.1 & 0.005 & 0.075 & 0.073 & 0.071 \\
1000 & 0.3 & 0.043 & 0.125 & 0.112 & 0.065 \\
1000 & 1.0 & 0.347 & 0.660 & 0.588 & 0.114 \\
1000 & 3.0 & 1.151 & 2.121 & 1.394 & 0.457 \\
\hline
5000 & 0.1 & 0.005 & 0.051 & 0.053 & 0.050 \\
5000 & 0.3 & 0.043 & 0.095 & 0.101 & 0.046 \\
5000 & 1.0 & 0.347 & 0.608 & 0.589 & 0.093 \\
5000 & 3.0 & 1.151 & 1.987 & 1.448 & 0.235 \\
\hline
20000 & 0.1 & 0.005 & 0.175 & 0.175 & 0.176 \\
20000 & 0.3 & 0.043 & 0.184 & 0.184 & 0.176 \\
20000 & 1.0 & 0.347 & 0.623 & 0.593 & 0.169 \\
20000 & 3.0 & 1.151 & 2.008 & 1.419 & 0.203 \\
\hline
\end{tabular}
\caption{$W_2(\hat p,p_0)$ for the three reconstructions.}
\label{tab:sensor}
\end{table}

\begin{figure}[H]
\centering
\includegraphics[width=\linewidth]{sensor_w2.png}
\caption{$W_2(\hat p,p_0)$ against $\sigma_{\mathrm{sensor}}$ for each $N$.}
\label{fig:sensor}
\end{figure}

As in 4(d), $W_2$ between two independent samples of $p_0$ is not zero, and with two modes it varies a lot from sample to sample
($0.57$, $0.13$ and $0.23$ for the three values of $N$ in our runs). Differences of this size are therefore not meaningful; the
conclusions below rely only on larger differences.
\begin{itemize}
\item \textbf{Small noise} ($\sigma\le0.3$): all three methods are equally good, since $Y$ is already close to $X_0$.
\item \textbf{Raw $Y$:} the error grows roughly in proportion to $\sigma$, to about $2.0$ at $\sigma=3$, because the noise widens both modes
and fills the gap between them.
\item \textbf{Shrinkage:} better than raw $Y$ at large $\sigma$ ($1.4$ against $2.0$ at $\sigma=3$), but still far off. Multiplying by $\lambda<1$
makes the spread smaller, but it also pulls the two modes towards each other and keeps them blurred. It is the best
\emph{linear} estimate of each $X_0$, not a reconstruction of the distribution.
\item \textbf{Reverse diffusion:} stays near the noise level for all $\sigma$ and $N$ (at most $0.46$, and $0.20$ at $N=20000$, $\sigma=3$). After
rescaling, the observations are exactly a sample of $p_{t^*}$, so the reverse process recovers $p_0$ up to sampling and
discretisation error. Note that it reconstructs the \emph{distribution}: each output is a new sample from $p_0$, not a denoised
version of the corresponding $X_0^{(i)}$.
\end{itemize}

\end{document}
